\begin{frame}{소실/폭발 갈라지는 세 곡선}
  \begin{columns}[T]
    \begin{column}{0.52\textwidth}
      \centering
      \includegraphics[width=\linewidth]{ch11_3_decay.png}
      \vskip0.3em
      {\footnotesize 승수 0.45 / 0.75 / 1.5의 그래디언트 곱이
      시퀀스 길이 $T$에 대해 지수적으로 소실 · 느린 소실 ·
      폭발로 갈라지는 곡선.}
    \end{column}
    \begin{column}{0.48\textwidth}
      \begin{block}{핵심 통찰}
        ``\textbf{승수의 밑}이 1보다 작으면 지수 소실, 크면 지수
        폭발''.
      \end{block}
      \vskip0.5em
      \begin{itemize}
        \item 초기화를 $W_{hh}$ 고유값이 1 근처에 오도록 정교하게
          맞춰야 한다는 11.2절의 ``자주 하는 실수''가 여기서 그
          \textbf{깊이}를 드러냄
        \item 그러나 어떤 초기화도 \textbf{지수}라는 형은 못 바꿈
        $\Rightarrow$ \textbf{구조}를 바꿔야 함
      \end{itemize}
    \end{column}
  \end{columns}
\end{frame}
